\section{Conclusion}
\label{sec:conclusion}

The reported reliability boundaries are empirical limits of supported gains within the tested English configurations. They distinguish expansion gains from evidence delivery and answer gains; unresolved gains do not establish absence of an effect. The hosted HotpotQA gain does not establish a content-independent effect or transfer to either pinned open generator. On news documents, standard HyDE underperforms matched Dense retrieval and all eight reader/length EM conditions. Better prompting improves the standard condition without establishing superiority over Dense.

The learned cross-encoder has supported retrieval gains over BM25 in the Wikipedia-oriented tiers and news documents. A random 300-question comparison supports its answer gains over BM25 with both readers, while other answer decisions differ across configurations. The final-input audit confirms character-stage support loss in the actual HotpotQA reader configuration, without additional token-stage loss. A matched 100-question budget change recovers additional support without establishing answer gains. These observations identify distinct measurement boundaries; they do not establish a causal explanation or human evidence sufficiency.

The findings motivate matched question-only baselines, effective-input checks, and direct answer comparisons when evaluating expansion or reranking. Named contrasts and family-specific paired decisions make those checks inspectable within the tested English configurations. Gold-selected controls improve both readers over question-only but do not establish a causal bottleneck or strict answer upper bound. The 2Wiki split-union corpus, unavailable external IRCoT, incomplete end-to-end costs and withdrawn independent human audit constrain generalization. New-sample validation does not establish a selection benefit when competing frozen rules select the same method. The study therefore provides neither a universal account of when expansion works nor a validated predictive model or method-selection rule; transport to untested configurations requires further evaluation.
